% LSPV-4 REPORT::WEIGHTPAD1::2026-10-01T20:46:00-07:00
% !TEX program = xelatex
\documentclass[11pt]{article}
\usepackage[margin=0.88in]{geometry}
\usepackage{amsmath,amssymb}
\begin{document}
\section*{LSPV-4 hostile audit: witness-weighting destroys normalized Search--G aggregates}
\textbf{STATUS: OPEN.} P versus NP remains open. This report proves representation-instability facts for Search--G/FNP witness aggregates and independently audits LSPV-3 QLEN1. It proves no unrestricted SAT lower bound, arbitrary-algorithm transfer, or nonrelativizing/nonalgebrizing hinge.

\paragraph{Predecessor preflight and CORPUS REVIEWED.}
The newest verified LSPV-3 artifact is \texttt{2026-10-01\_2035\_LSPV-3\_QLEN1.tex}; it is substantive, correctly headed, sentinelled after \texttt{\string\end\{document\}}, and readable from GitHub. I surveyed the current repository tree; AGENTS.md, CODEX.md, the pnp-research skill, root main.tex, TASK-001 through TASK-007, newest verified MAIN 1802 and relevant MAIN 1558, RANKTRANS1, SUPPORTK1, QLEN1, NOGUARD1, ACQ1, REREAD1, PCLOSE1-R1, RANDCLOSE1, REACQ1, PADINV1, ISOCAN1, governance, HEC/VSC-HEC/G-PPC-HEC/CRS-PPC/CCSPC/RCCN, and the FP/FNP Search--G spine. Repository search covered STORE/REACQUIRE, counting/parity, Search--G, ACQ1, arbitrary-algorithm transfer, and duplicate causal claims.

\paragraph{Primary-literature preflight and obstacles.}
The immediate Search--G obstacle is whether PADINV1's surviving normalized witness fraction is representation-stable after raw count/parity were retired. The reconstruction obstacle is whether QLEN1's improved SAT-oracle substitution is valid in its declared explicit-query-writing model. Valiant's counting framework separates counting from ordinary NP search; Valiant--Vazirani shows multiplicity can be altered by randomized isolation; Toda shows the strength of PP/counting relative to PH. None forces an arbitrary FNP selector to compute a witness aggregate. Aaronson--Wigderson remains the barrier reference.

\paragraph{PROVED/EXACT: selector-preserving weighted duplication.}
Let $R(x,y)$ be polynomial-time decidable and polynomially balanced, and let $h(x,y)\in\{0,1\}$ be polynomial-time. Fix polynomial $p(n)\ge1$. Define
\[
R^{+}_{p,h}(x,(y,z))\iff R(x,y)\wedge(h(x,y)=1\ \vee\ z=0^{p(|x|)}),
\quad |z|=p(|x|).
\]
Writing
\[
A=|\{y:R(x,y),h(x,y)=1\}|,\qquad B=|\{y:R(x,y),h(x,y)=0\}|,
\]
gives
\[
|W_{R^+}(x)|=2^{p(n)}A+B,\qquad
\boxed{\mu^+_{p,h}(x)=\frac{2^{p(n)}A}{2^{p(n)}A+B}}.
\]
If $A>0$, the normalized $h$-mass tends exponentially to one as $p$ grows. The search problems are polynomially selector-equivalent: a selector $S$ for $R$ yields $(S(x),0^{p(n)})$ for $R^+$, and projection of any $R^+$ selector yields an $R$ selector.

Similarly,
\[
R^{-}_{p,h}(x,(y,z))\iff R(x,y)\wedge(h(x,y)=0\ \vee\ z=0^{p(|x|)})
\]
has
\[
\boxed{\mu^-_{p,h}(x)=\frac{A}{A+2^{p(n)}B}},
\]
which tends exponentially to zero when $B>0$, again with polynomial selector equivalence.

\paragraph{Consequence for PADINV1 and Search--G.}
PADINV1 correctly proved that uniform COPY preserves a padding-independent normalized witness fraction. WEIGHTPAD1 shows that witness-dependent duplication destroys that invariant while preserving the search problem up to polynomial selector equivalence:
\[
\boxed{\text{normalized witness fraction is not a generic semantic invariant of FNP search}.}
\]
Generic normalized witness mass as a Search--G hardness resource is therefore \textbf{FALSIFIED/RETIRED}. This does not alter a fixed source-sanctioned relation. A surviving aggregate must justify why its encoding/multiplicities are canonical and must include a selector-to-aggregate forcing theorem.

\paragraph{Counting/isolation guard.}
Nothing here proves $\#P=FP$, $PP=P$, or $\oplus P=P$ under $P=NP$. Valiant, Valiant--Vazirani, and Toda remain KNOWN THEOREMS about counting/isolation. The hostile point is narrower: raw count, parity, normalized mass, or uniqueness attached to an arbitrary FNP presentation does not force every selector to pay for that aggregate.

\paragraph{Independent red-team audit of QLEN1.}
QLEN1's local query-writing lemma is correct in its declared oracle-Turing-machine model. If writing each query symbol costs one ordinary step and branch time is $T$, then $\sum_i s_i\le T$. If $A(s)\le s^{\kappa+o(1)}$, $\kappa\ge1$, then
\[
\boxed{\sum_i A(s_i)\le T S_{\max}^{\kappa-1+o(1)}}.
\]
This is amortized accounting, not a lower bound, and does not apply unchanged to a unit-cost interface naming long queries. The local inequality is independently confirmed as \textbf{PROVED/EXACT in the declared model}. The specialized vMS/SGL exponent ledger remains \textbf{CONDITIONAL} on source certification of $S_{\max}$, amplification, compiler exponents, and SAME\_REALIZATION model matching. TC4 remains OPEN.

\paragraph{SU/vMS, acquisition, reconstruction, and PQ/cross-domain ledger.}
No structural claim is strengthened. RANKTRANS1/SUPPORTK1 sharpen fixed-code algebra but exact SU09/vMS invocation, legal bit-label realization, stacked Schur/common-kernel rank, and extension across legal lines remain OPEN. ACQ1 remains CONDITIONAL on the exact legal-state evaluator and SAT compilation exponent. Sparse-$Q$ normalization remains OPEN. Fixed-code Schur products, quadratic kernels, Boolean polarization/exterior identities, typed puncturing/shortening, cyclic convolution, Gaussian elimination, and QLEN1's explicit query-writing inequality are \textbf{EXACT OBJECT/THEOREM TRANSFERS}. Syzygies, Betti invariants, minimum-rank relations, shortening filtrations, collision/special-soundness extraction, and LWE-style black-box extraction are \textbf{ANALOGY/CANDIDATE}. Cryptographic hardness/random oracles, VLSI locality, distributed cuts, online no-reread bounds, control/observability, topology, and proof-system-specific lower bounds are \textbf{FAILED TRANSFERS} to unrestricted static SAT runtime.

\paragraph{Hostile controls and barrier audit.}
COPY implements weighted duplication; relabeling defeats naive lexical semantics. DEAD/DELAY inflate presentation without semantic work. INTERP/RECOMP, static reread, memoization, pointer jumping, and group compression remain bypasses. Gaussian/Tseitin, Horn/2-SAT, bounded treewidth, independent products, padding, semantic oracles, and algebraic oracles remain controls. Polynomial latent branching is not NEXP. WEIGHTPAD1 and QLEN1 relativize; the active structural branch remains algebraic. No nonrelativizing/nonalgebrizing hinge is obtained.

\paragraph{DUPLICATION / CONTRADICTION AUDIT.}
PADINV1 retired raw count/parity and observed uniform-padding invariance of normalized fractions. ISOCAN1 retired unique/canonical witness reacquisition. WEIGHTPAD1 is new: nonuniform witness-dependent duplication destroys normalized-fraction invariance while preserving selector equivalence. No contradiction with RCCN STORE-or-REACQUIRE occurs. QLEN1 does not contradict SGLCOST1; it sharpens substitution under the additional explicit-query-writing premise.

\paragraph{STALE OR OVERSTATED MAIN CLAIMS; WHAT CHANGED RELATIVE TO ROOT MAIN.}
Any root-main reading that treats all-witness normalized mass as a generic escape merely because uniform COPY preserves it is stale. Raw multiplicity, parity, lexical canonicity, uniqueness, and normalized witness mass are all insufficient generic selector-hardness resources. QLEN1's local amortized substitution inequality may now be promoted after independent audit, while its full SGL reconstruction ledger remains CONDITIONAL.

\paragraph{ADOPTED / REJECTED / RETIRED PRIOR COMPONENTS.}
Adopted: source-sanctioned Hadamard typing, fixed-code exact algebra, charged ACQ1, RCCN STORE-or-REACQUIRE in its declared model, REACQ1, PADINV1, ISOCAN1, and QLEN1's local lemma. Rejected as standalone hardness: raw count, parity, normalized fractions, uniqueness, lexical canonicity, delayed challenge, and causal provenance. Retired: generic promotion of these presentation-level witness aggregates to a superpolynomial theorem for arbitrary FP/FNP selectors.

\paragraph{P=NP-compatible components without claiming P=NP.}
Static reread, deterministic selector rerun, canonical witness recovery, polynomial adaptive SAT calls, bounded-error polynomial causal computation, stagewise nondeterminism, memoization, pointer/group compression, selector-preserving weighted duplication, and ACQ1-style NP search are compatible with the counterfactual. They are hostile controls, not evidence for $P=NP$.

\paragraph{Strict status and next target.}
\textbf{PROVED/EXACT:} weighted-duplication identities; polynomial selector equivalence of $R$ with $R^+_{p,h}$ and $R^-_{p,h}$; representation instability of normalized witness mass; QLEN1 local query-writing substitution in its declared model.
\textbf{KNOWN THEOREM:} Valiant counting, Valiant--Vazirani isolation, Toda counting/PH, Aaronson--Wigderson algebrization.
\textbf{CONDITIONAL:} ACQ1 and specialized QLEN1/vMS reconstruction ledger.
\textbf{FALSIFIED/RETIRED:} raw count/parity, uniqueness/canonicality, normalized witness fraction as generic FNP selector-hardness resources; generic STORE-or-REACQUIRE promotion; prior polynomial-depth causal/randomness escapes.
\textbf{OPEN:} exact SU09/vMS invocation; stacked Schur/common kernel; sparse-$Q$ normalization; exact ACQ1 exponent; reconstruction source parameters; arbitrary-algorithm transfer/TC4; a representation-invariant selector-forced resource; nonrelativizing/nonalgebrizing hinge; final contradiction.

The next falsification target is the strongest surviving Search--G semantic aggregate: require invariance under efficient witness relabeling and witness-dependent selector-preserving duplication, then attack the exact theorem claimed to force every selector to acquire it.

\begin{thebibliography}{4}
\bibitem{Valiant79} L. G. Valiant, \emph{The Complexity of Enumeration and Reliability Problems}, SIAM J. Comput. 8(3):410--421, 1979.
\bibitem{VV86} L. G. Valiant and V. V. Vazirani, \emph{NP is as Easy as Detecting Unique Solutions}, Theor. Comput. Sci. 47:85--93, 1986.
\bibitem{Toda91} S. Toda, \emph{PP is as Hard as the Polynomial-Time Hierarchy}, SIAM J. Comput. 20(5):865--877, 1991.
\bibitem{AW09} S. Aaronson and A. Wigderson, \emph{Algebrization: A New Barrier in Complexity Theory}, ACM Trans. Comput. Theory 1(1):2:1--2:54, 2009.
\end{thebibliography}
\end{document}
% END_OF_VERIFIED_REPORT::LSPV-4::WEIGHTPAD1